\documentclass{ltjsbook}
\usepackage{docmute} 
\usepackage{../../myStyle} 

\begin{document}
\chapter{プログラミングの基礎}
\label{x16_Programming1}
この授業は，RやStanといったプログラミング言語を扱います。プログラミングはやったことがない，という人でも理解できるように，ゆっくり話を進めていきます。
プログラミングのいいところは，人間相手とのやり取りと違って，必ず理屈があって振る舞いが生じている点です。対話を通じて根気よく付き合えば，必ず問題は解決できます。また昨今では，生成AIがプログラミングをアシストしてくれるようになりました。この授業はTAがついていますから，わからないところがあれば，随時TAに尋ねてくだされば結構です。それでも人間相手の質問ですから，「こんなことを聞いたらバカだと思われるんじゃないだろうか」「声をかけるのが恥ずかしい」ということがあるかもしれません。でも生成AIとのやり取りは，いつでも気軽に行えます。また決してこちらを蔑んだり恨んだりすることはありません(というかあっても怖くありません！)。同じ質問を何度しても，生成AIは一向に意に介さないので，積極的に利用してください。ただし生成AIは正しい答えを教えてくれるとは限りませんので，\textbf{必ず自分が確認できるような質問をしてください}。何かわからないけど，バグが出ないからいいや，と思っていると，自分がなんの成長もしないのはもちろん，大変な問題を引き起こすかもしれませんので。

\section{パソコンの基礎}

21世紀に生きる私たちは身の回りをコンピュータに囲まれて生きています\footnote{私は1976年生まれですが，生まれた頃は周りにコンピュータなんかありませんでした。小学生の頃にマイコン(マイクロコンピュータ，小さなコンピュータという意味でもありますが，My Computer，私のコンピュータという意味でもあります。つまり個人単位でコンピュータが使えるようになった，というだけでも大きな出来事だったのです。)という言葉が出てきて，なんかかっこいいなと思った記憶があります。私が10歳になったころ，ビデオゲーム(テレビゲームでやるゲームが家庭でできるようになり，それをこのように呼びました。)が身の回りに出てきました。ファミリーコンピュータ，とくに「スーパーマリオブラザーズ」によって日本中の子供たちが熱狂したのが11歳の頃，「ドラゴンクエスト」によって社会問題になったのが12歳の頃になります。ともかくこの頃は，コンピュータといってもゲーム機のような扱いでした。}。それは携帯電話の形をしていたり，ノートパソコン，デスクトップパソコンの形をしていたりします。また時々テレビなどで報道されますが，気象予報や飛沫がどのように飛び散るかをシミュレーションする大型計算機「富嶽」などもコンピュータですね。

みなさんはスマートフォン，タブレット，ゲーム機には慣れ親しんでいると思います。パソコンも「GIGAたん」などで触れていると思います。パソコンでは文書，プレゼン，図表を作ったり，絵や音楽を作ったりしているでしょう。このようにパソコンを\textbf{使う}ことには慣れていても，パソコンに\textbf{命令する}という感覚はないかもしれません。

プログラミングとは，計算機に計算をさせる指示書を書くことです。パソコンをはじめとするコンピュータ機器は，その形状はさまざまですが，いずれも電子計算機であり，なんらかの形で計算をさせていることになります。ゲームで遊ぶことも，動画を見ることも，文書を作ることも，すべて電子的な計算の結果なのです。プログラミングをする上では，基本的な計算機の仕組みを知っておいても損はないでしょう。

電子計算機には次の5つの装置があります。

\begin{description}
	\item[入力装置] キーボードやマウス，タッチパネルなどを使って情報を取り込むデバイス(装置)
	\item[出力装置] モニタやプリンタ，タッチパネルなどを経由して情報を出力するデバイス
	\item[演算装置] プログラムの命令に従って計算処理(四則演算や論理演算)をする装置。一般に電子計算機の中央で一括して処理するので，Central Processing Unit(CPU)と呼ばれるものです。
	\item[制御装置] 演算結果に従って他の装置に指示を出す装置のこと。演算装置とまとめてCPUに実装されています。
	\item[記憶装置] 計算結果などの情報をいったん保持しておく装置のこと。コンピュータの内部にあって一時的な計算に使われる一次記憶装置，ハードディスクドライブ(Hard Disk Drive,HDD)やソリッドステートドライブ(Solid State Drive,SSD)などの二次記憶装置など。
\end{description}

最後の記憶装置については，一次記憶装置，二次記憶装置と種類が分かれていますがこの区別は簡単で，電源を落とした時に記憶が消えてしまうのが一次記憶装置，電源を落としても記憶が消えないのが二次記憶装置です。たとえば$12+38=$という計算をするとき，頭の中で「えーっと一の位が2と8だから10になって繰り上がるから・・・」と考えてから，ノートに$=50$という答えを書くと思います。次の問題に進むと，先ほどの「1繰り上がるから・・・」という情報は忘れてますよね。でもノートに書いた$12+38=50$というのは残っています。このノートがいわば二次記憶装置であり，頭の中で一時的に保持していた情報が一次記憶装置ということになります。一次記憶装置はRAM(Random Access Memory)とも呼ばれます\footnote{実はこのように人間を1つのコンピュータに喩えて，そこではどのように計算がされているのか，人間の記憶装置や演算装置，入出力装置の特徴はどうなっているのか，というのを研究するのも心理学の仕事です。記憶や演算，制御については認知心理学や学習心理学，入出力については知覚心理学や生理心理学が専門的に扱っています。そういう意味でもコンピュータの登場は心理学に大きな影響を与えているのですが，それはまた別の講釈で。}。

ところで，コンピュータがやっているのは計算だけです。私はこの資料をPCに向かって書いており，キーボード(入力装置)を叩きながら，画面(出力装置)を見て文字を連ねています。これも「キーが押されたら文字を表示させ，その文字列を記録する」という処理を機械が淡々とこなしているに過ぎません。あるいはマウスやトラックパッドで，アイコンを指し示し，カチリと押す\footnote{押すときの音から，この操作をクリックclickと言います。2回続けて押すことをダブルクリックと言います。}ことで選択し，押し込んだまま移動させ(ドラッグ)，離すことで別の場所に置いたりします(ドロップ)。トラックパッドの場合は2本指で同時に押したりしますし，タッチパネルの場合は二本指を広げたり(ピンチアウト)，逆に二本指を狭めたり(ピンチイン)，3本以上の指でファサーっと触って場所を広げたりします。たとえば「ファイルを掴んでゴミ箱に捨てる(削除する)」というのが我々にとっての操作ですが，コンピュータの内部では実はこんなことをしていません。ファイルは(二次)記憶装置に書き込まれた情報です。記憶装置は原稿用紙のように小さなマス目がたくさんあって，ファイルとはそのマス目のXXX番目からYYY番目までの情報，ということです。このファイルを削除するというのは，記憶装置のある場所(アドレス)に「削除されたものなので画面に表示しない」という情報を書き込む，という操作をしているだけです。じゃあなぜ私たちは「ゴミ箱にドラッグ＆ドロップ」なんてするのでしょう？ それはそのほうがわかりやすいからですよね。「ファイルを削除する」というのは，「メモリアドレスのXXX番地に別の情報を書き込む」という操作だと言われてもピンとこないので，コンピュータが人間にとってわかりやすい表現をして見せてくれているのです。このユーザにとってわかりやすい幻を見せてその気にさせてくれるというデザインのことを，ユーザーイリュージョンと言います。ともかくこういう「画面で見ながら操作する」ことをグラフィカル・ユーザ・インターフェイス(Graphical User Interface,GUI)といいますが，これのおかげでコンピュータの操作は随分楽になっています\footnote{実は人間の意識もこのユーザーイリュージョンのようなもので，実際の体の動かし方や感覚情報の受け止め方，処理の仕方は別ですよね。すべての情報を意識に上げるのではなく，「私はXXXをしている」と身体がそれっぽい幻想を投影して見せてくれているのが意識の正体ではないか，という議論があります。興味のある人は\textcite{UserIlluision2002}を読んでみてください。}。

プログラミングとはこうした装置を使う指示書を書くことでもあります。この授業で扱うのは主に数値計算ですから，逐一「この装置にこの動きを命令」といったところまで書き込む必要はありません。そのあたりのことは，言語環境がやってくれますのできにする必要はないのですが，根本的なところで何が動いているかはイメージしておいた方が良いでしょう。


\section{プログラミングの基礎}
あらためて，プログラミングとは，コンピュータに計算をさせるその仕様書を作ること，だと思ってください。お料理のレシピのように，計算レシピを書くわけですね。ただ相手は計算機ですので，言葉の端々まで正確に伝えないと理解してくれません。よく「思った通りに動いてくれない！」と不満を訴える初心者がいますが，それは当然で，プログラムは\textbf{思った通りに動くのではなく，書いた通りに動く}のです。思った通りに動かないのであれば，それは仕様書・レシピの方に間違いがあります。

今の第一原則に加えて，プログラミングを進める上での注意点をあげておきます。簡単なことだと思うかもしれませんが，基本的なルールをしっかり守ることが上達への近道です。
\begin{enumerate}
	\item プログラムは思った通りには動かない。書いた通りに動く。大文字，小文字，スペルの違いに注意する。
	\item 書き間違えないための工夫は美しさ。綺麗に書くことが大事。
	\item 一言一句すべてに意味がある。ただの写経ではなく，意味を考えながら書く。
	\item 「遊び心」をもって！ここを変えたらどうなるか，を少しずつ「やってみる」が大事
	\item 変更は少しずつ。一気に変えるとどこが変わったかわからなくなるから。
\end{enumerate}

\subsection{綺麗に書こう}
綺麗に書く，というのはコードの可読性を上げるために必要です。綺麗な書き方として，\textcite{Matsuura201610}は次の5点をあげています。

\begin{itemize}
	\item インデントは必ずする
	\item データを表す変数の先頭の文字は大文字。パラメータを表す変数の先頭の文字は小文字。
	\item 各ブロックの間は１行あける
	\item camelCaseではなく\verb|snake_case|で。
	\item 「~」や「＝」の前後は1スペースあける
\end{itemize}

1つめのインデントとは，字下げのことです。行の頭を少し凹ませることで，違う人まとまりであることを明示するのです。たとえば次の2つのコード\footnote{このコードは確率的プログラミング言語Stanの書き方で，この科目の後半で学ぶないようですから，今その中身までは理解しなくても結構です。}は同じ働きをしますが，\ref{code:16_02}の方が見やすいと思いませんか？
\begin{lstlisting}[language=c++,label={code:16_01},caption={インデントのないコード}]
data {
int<lower=0> N;
array[N] y;}
parameters {
real mu;
real<lower=0> sigma;}
model {
for(i in 1:N)
y[i]~normal(mu, sigma);}
}
\end{lstlisting}


\begin{lstlisting}[language=c++,label={code:16_02},caption={インデントのあるコード}]
data {
  int<lower=0> N;
  array[N] y;
}

parameters {
  real mu;
  real<lower=0> sigma;
}

model {
    for( i in 1:N ){
        y[i] ~ normal(mu, sigma);
    }
}
\end{lstlisting}

インデントのある\ref{code:16_02}は，\texttt{data}というブロック(\{ と\} で囲まれている領域)の中に，2つの行が入っている，ということが明確にわかります。また\texttt{model}というブロックの中には，\texttt{for}から始まる分がありますが，これも複数行に渡るブロックを構成する文なので，その中身はインデント(字下げ)されています。このように，インデントすることでどこの列がどこまでブロックを組んでいるのかがわかりやすくなるのです。ちなみにこのインデントを作るのはTABキーを使います。TABキーは普段，どう使うのかわからないものだと思われがちですが，このインデントをしてくれるためのものです\footnote{ここには流派があって，TABでインデントする派とスペースキーを4回または8回押してインデントする派がいます。どちらでも働きは同じなのですが，個人的には数回の字下げを1回のキーで行ってくれるTABのほうが良いと思っています。}。コードエディタによっては，カッコで括ったときに自動的に閉じるカッコを用意し，また改行ごとに字下げ位置を合わせてくれるものがあります。これらの機能を使って是非わかりやすいコードを書いてください。

\textcite{Matsuura201610}の指摘の2,3番目についてはお好みで，4番目もそれほど強い・広く浸透した決まりではありません\footnote{ちなみにキャメルケースCamel Case，スネークケースSnake Caseとは変数名のつけ方のルールです。Rなどプログラミング言語ではあらゆるものに名前をつけて管理しますが，名前のつけ方は任意です。命名はわかりやすい方がいいですが，長いものになると区切りを入れたくなるかもしれません。さてCamelとはラクダの意味で，ラクダのコブのように区切りのところだけ大文字にする，という記法です。Snakeは蛇のことで，区切りのところにアンダースコア\verb|_|を使うというものです。たとえば野球チーム，という変数名をつけたい時に，キャメルケースのやり方だと\texttt{BaseballTeam}のように書きますし，スネークケースのやり方ですと\verb|Baseball_team|のような書き方になります。ちなみにRは日本語での命名も許しますが，そのほかの言語では一般的ではありませんし，何より全角と半角を切り替える時のミスやエラーが多くなるので，半角英数字での命名をすべきという点はどちらでも共通です。}。ただし，5番目のスペースの前後に少し空白を取るのは，見やすさのためにも是非実行してもらいたいところです。

これらの書き方は，\keytermL{可読性}{かどくせい}を上げるためのものです。プログラムは仕様書・レシピですから，あとで読み直した時やほかの人が読むときも，意味がわかるようにしておいた方が良いのです。自分だけわかれば良い，と思うかもしれませんが，その自分ですら何があるのかわからなくなってしまう可能性があります。そのような読みにくさはすぐにバグにつながりますから，綺麗に書くことを常々心がけておいて欲しいのです。

\subsection{意味を考えて書こう}
プログラミング言語は，ほとんど英単語のようなものです。プログラミング上の都合から，短い言葉に省略されていたり，アンダースコアやピリオドでつながって書いてあったりしますが，比較的わかりやすい言葉であることに違いはありません。

プログラムを習得し上達するためにすべきことは，最初は写経と呼ばれる，テキストや指示にしたがって書き写すことです。もちろんネット上にあるサンプルコードなどをコピー＆ペーストしても良いのですが，自分用の分析コードを書くためには手を入れる必要があったりします。ですから，最初はなるべく自分でキーボードを叩いて，コピー＆ペーストではなく入力する経験を積んでください。もちろん間違えてしまうことが多いでしょうが，転ばずに歩けるようになった人が一人もいないように\footnote{お釈迦様は除く。}，ミスを犯すことでどういうミスだったかを考え，自覚することではじめて，すこしずつですがbミスが減っていきます。失敗する経験も時には必要なのです。

ということで，時折自分なりにコードを書いてもらうのですが，その時に「ただの記号列」と思わないでください。すでに書いたように，英語あるいはそれを省略した表現になっているので，見慣れないものかもしれませんが必ず意味があります。
プログラムを始める最初の頃は，ミススペルが多く，あちこちでエラーが出て気持ちが萎えることもあるかもしれません。エラーが出る時は目を凝らして，どこに問題があるかを考えて修正することになります\footnote{ちなみにプログラミング歴30年以上の私でも，簡単な英単語，たとえばlibraryのスペルを間違えるなんてことはしょっちゅうです。}。ここで難しいのが，$x$や$s$などの大文字と小文字の区別がつきにくいもの，$1$と$l$のように違いがわかりにくいものがあるということです。
私は職業柄，多くの人に教え，多くのバグを見つけてきましたが，その中でも特別見つけるのに苦労したのが，\texttt{norma1}というミススペルでした。みなさん，これのどこがミススペルかわかりますよね？でもこれ，書く時に「正規分布のことだな」と思っていれば，そもそも入力時にそんな入れ方はしないと思うのです。ただの字だ，と心を無にしてしまうと，後々見つけにくいエラーを作ることにもなりますから，この文字・この名前・この関数はどういう意味だろう，と少し考えながら取り組んでみてください。

\subsection{遊び心が大事}
プログラミングを進める上では，遊び心が大事です。うまく動くコードがかけたら，もうおかしなことになりたくない，と手をつけないままにしてしまう人がいます。

ところで，今まで教えてきた経験からいって，プログラミングが上達する人は，うまくコードが書けたらすぐに「ここをこう変えたらどうなるんだろう」と試してみる人だと思います。最初はプログラムの意味がまったくわからないので，どこをどういじっていいのか見当もつかないかもしれません。しかしたとえば\texttt{blue}という文字が出てきたら，これは色の青のことじゃないか，と思いますよね。ではその\texttt{blue}を\texttt{red}に変えたらどうなるんだろう？やってみよう！となるような，そういう遊び心が大事なのです。

もちろん変えてしまったことでエラーになって，動かなくしてしまうことがあるかもしれません。その場合は，元に戻して元通り動くかどうか，さらに確認すれば良いのです。数字や文字を少しいじっただけで，PCが爆発してしまうようなことはありませんから，ちょっと遊び心を出して，どうなるのかな？と思う気持ちを大事にしてください。

\subsection{変える時は少しずつ}
プログラムを色々いじって遊ぶことが成長につながる，という話をしました。ただし遊び方には気をつけましょう。まずうまくいくコードができれば，それは保存しておいて，また同じ内容のものを別名で保存し，「遊んでいい方」「壊れてもいい方」を作って，そちらで色々変えて遊ぶといいでしょう。最悪なことがあっても，うまくいくバージョンは常に残っているわけですから。

そして色々試す時のコツは，一箇所ずつ変更していくことです。たとえば\texttt{blue}を\texttt{red}に変え，\texttt{line}を\texttt{box}に変え・・・と複数箇所を同時に変更して，うまくいかなかった場合，どちらが原因になっているのか把握できませんよね。これは心理学実験でいうところの交絡です。操作が交絡してしまって，原因が特定できなくなるのです。

また，実行するときも1行ずつやりましょう。とくにRは一問一答型，つまり1つの命令について1つの反応を随時返してきます。複数行をまとめて実行することもできますが，まとめて実行している途中でエラーが出ていて，その後の計算がすべて空回りしているということも少なくありません。エラーがどこで生じているのか，しっかり特定することが重要です。そのためにも焦らず，1つずつ確実にできることを積み重ねていくという姿勢が重要です。

統計分析も複雑で細部まで設定し尽くしたモデルを作り上げていくことができますが，いきなり全体が完成するのではなく，小さなピースの積み重ねなのです\footnote{Rやプログラミングで統計を学ぶ意義はここにあります。GUIで操作できる統計パッケージも少なくありませんが，それらは画面が出てきた時にデフォルトで設定されている値があるのがほとんどで，自分が何をやっているのか細部まで気づかないまま実行できてしまうのです。そうすると当然，エラーが出たり，うまくいっても誤用している，ということにもなりかねません。自分で理解できていない技術に振り回されないようにするために，しっかりとわかるピースを積み重ねていく必要があるのです。}。

\section{いくつかのプログラミング言語}

前置きが長くなりました。それではプログラミングを始めていきましょう。
プログラミングには，プログラムを計算機への命令文として伝達する「言語」を利用することになります。
この授業ではRという言語を使います。R言語は統計計算に特化した言語であり，心理学はもちろん統計に関する領域では多岐にわたって利用されています。比較的初心者に優しい言語設計になっており，また導入当初から豊富な関数，可視化(グラフ化)機能をもっています。統計解析を含んだ学術論文でもRを使ったものが多く，SPSSやSAS，Stataといった統計アプリケーションよりも広く使われています。その理由の一つは，Rがフリーソフトウェアであるということです。マイクロソフトやIBMといった企業が提供するソフトウェア(プロプライエタリといいます)は，企業秘密の名目で，アプリの中で何がどのように計算されているかが隠されています。計算結果に間違いはない，と信じたいところですが，隠れている部分がある以上完全には信用できません。Rはオープンソースのフリーソフトウェアです。すなわち，全ての計算過程を公開しており，お金をとって保証やサービスに代えるということをしていません。科学という営みにおいては，公明正大であることが特に重要ですから，この点は強調してもしすぎるということはないでしょう。

ところで，最近流行のAI，人工知能，機械学習といった領域では，RよりもPythonの方がよく利用されています。PythonもRと同様にオープンソースのフリーソフトウェアですし，優れた計算関数がパッケージとして提供されています。PythonでもRと同様のことができますが，Rよりもプログラミングにおいて厳格さを要求すること，複数のバージョンが乱立・併用して使われていることなどが理由でこの授業での採用をみおくりました。統計計算という意味では，Juliaという言語もありますので，興味がある人はさまざまな言語を使ってみてください。

ところで，R，Python，Juliaといった言語は\keytermL{インタプリタ}{インタプリタ}型の言語です。Interpret(翻訳)，という言葉からわかるように，機械に命令をするときに逐一命令文を翻訳して伝えます。「足し算をしろ(命令)」「したよ(結果)」「掛け算をしろ(命令)」「したよ(結果)」というように，一問一答型でプログラムが進んでいきますから，自分が今何を命令して，どのような答えが出たかが明確です。

これとは別に，\keytermL{コンパイラ}{コンパイラ}型言語というのもあります。C，Javaなどが代表例ですし，この授業の後半で利用するStanという言語もコンパイラ型です。この方式は，一連の作業工程をかいたプログラムを，まるごと翻訳してから計算を始めます。この翻訳のことをCompileする，といいます。一連の作業工程のどこかにミスがあれば--それがスペルミスのような単純なものであっても--翻訳の時にエラーが出ます。エラーには「XX行目でエラーが出たよ」というメッセージが表示されますが，必ずしもその行に問題があるとは限りません。その命令文の前の方で問題が生じていて，発見したのがたまたまそこになっただけという可能性もあります。さらにいえば，コンパイル時はエラーが出ず，計算が実行されてから，計算結果が間違っているとか，変な結果が表示されるということもあります。その時は元の設計書のどこに問題があるか，目を皿のようにして探し回る必要があります。インタプリタに比べて，どうにも不便が多いようにも思えますが，コンパイラの利点はその計算速度です。インタプリタは逐一のやり取りになるので，どうしても計算が遅くなりがちです。一気に計算させたい時，まとまった計算単位がある時はコンパイラ型の方が有利なのです。

\section{R言語の基本的な挙動}
前置きが長くなりました，と言ってからさらに前置きをしてしまいました。それでは早速Rを使っていきましょう。
Rはそのままでも使えますが，RStudioを使って利用した方が便利です。すでに一年時に，RStudioの基本的な使い方については触れていると思います。

RStudioをひらいて，この授業用に新しく\textbf{プロジェクト}を準備しましょう。プロジェクトを設定することで，ファイルの管理がやりやすくなります。ソースコードに簡単な四則演算を書いて実行し，また簡単な関数(\verb|sqrt|や\verb|help|など)を実行してみてください。問題なく利用できるでしょうか。以下に第一回目の課題を記しますので，早速やっていきましょう。

もしうまくできない場合は，教員やTA，周りの人を頼ってください。あるいは，生成AIに質問してもいいでしょう。
生成AIはうまく利用すると，Rのコードも書いてくれたりします。簡単な問題であれば，結構正しく動くコードを返してくれます。
ただし，生成AIのいうことを丸呑みにせず，必ず自分で試しながら「意図した通りのことができているか」を確認するようにしましょう。
\section{課題}
\subsection{RとRStudioの基礎}
\begin{enumerate}
	\item この授業のために，新しいプロジェクトを作成してください。プロジェクトは新しいフォルダでも，既存のフォルダでも構いません。
	\item プロジェクトが開いた状態のとき，RStudioのウィンドウ・タブのどこかに「プロジェクト名」が表示されているはずです。確認してください。
	\item 新しいRスクリプトファイルを開き，空白のままで結構ですからファイル名をつけて保存してください。
	\item RStudioを終了あるいは最小化させ，OSのエクスプローラ/Finderから，プロジェクトフォルダに移動してください。先ほど作ったファイルが保存されていることを確認してください。
	\item プロジェクトフォルダには，プロジェクト名+.projというファイルが存在するはずです。これを開いて，RStudioのプロジェクトを開いてください。
	\item RStudioのFile > Close Projectからプロジェクトを閉じてください。画面の細部でどこが変わったか，確認してください。
	\item RStudioを終了し，再びRStudioを起動してください。起動の方法はプロジェクトファイルからでも，アプリケーションの起動でも構いません。起動後に，プロジェクトを開いてください(あるいはプロジェクトが開かれていることを確認してください。)。
	\item Rを起動し，新しいスクリプトファイルを作成してください。そのファイル内で，2つの整数を宣言し，足し算を行い，結果をコンソールに表示してください。
	\item スクリプトに次の計算を書き，実行してください。
	\begin{enumerate}
		\item $\frac{5}{6} + \frac{1}{3}$
		\item $9.6 \div 4$
		\item $2.3 + \frac{1}{2}$
		\item $3\times (2.2 + \frac{4}{5})$
		\item $(-2)^4$
		\item $2\sqrt{2} \times \sqrt{3}$
		\item $2\log_e 25$
	\end{enumerate}
\end{enumerate}

\end{document}
